\documentclass[10pt,a4paper]{amsart}
\usepackage[margin=19mm]{geometry}
\usepackage{amsmath,amssymb,mathtools}
\usepackage[scr=rsfs,cal=euler]{mathalfa}
\usepackage{xfrac}
\usepackage[unicode,hidelinks,bookmarks=false]{hyperref}
\newcommand{\ie}{i.e.\ }
\newcommand{\cf}{cf.\ }
\newcommand{\resp}{respectively\ }
\newcommand{\Subsection}{\subsection}
\providecommand{\nobreakdash}{\nobreak\hbox{-}}
\setlength{\emergencystretch}{2em}
\multlinegap0pt
\usepackage{bm}
\usepackage{bbm}
\usepackage{dsfont}
\usepackage{xspace}
\usepackage{cancel}
\usepackage{mathtools}
\usepackage{amsthm}
\makeatletter
\@ifundefined{theorem}{\newtheorem{theorem}{Theorem}}{}
\@ifundefined{claim}{\newtheorem{claim}[theorem]{Claim}}{}
\@ifundefined{lemma}{\newtheorem{lemma}[theorem]{Lemma}}{}
\makeatother
\title[Subdivision-free graphs with the maximum spectral radius]{Subdivision-free graphs with the maximum spectral radius}
\author{Wanting Sun, Guanghui Wang, Pingchuan Yang}
\date{Source: arXiv:2507.04257v1 (CC BY 4.0)}
\begin{document}
\maketitle

\noindent\emph{Source and licence.} Subdivision-free graphs with the maximum spectral radius. Version arXiv:2507.04257v1 (CC BY 4.0), distributed under the Creative Commons Attribution 4.0 International licence, as recorded in the arXiv licence metadata for this exact version and shown on the abstract page https://arxiv.org/abs/2507.04257v1. Full rights notice: licenses/arxiv-2507.04257-CC-BY-4.0.txt.\par\smallskip

\noindent\textit{Editorial scope.} Counted unit: the Claim whose source label is \texttt{lem3.5} in the article (source0.tex lines 573--575, source label \texttt{lem3.5}), delivered together with the article's own complete original proof of it (source0.tex lines 576--648, one \texttt{proof} environment of 73 lines). No mathematics is written, repaired or re-derived: statement and proof are the article's own text line for line. Required context retained uncounted: lem3.1 (lines 288--291); lem3.2 (lines 306--308); lem3.4 (lines 540--542). Every definition, display and label of those lines is retained unchanged. Omitted from this package and not redistributed: 1--287, 292--305, 309--539, 543--572, 649--1098. Nothing inside a retained excerpt is omitted or altered: every retained body is delivered exactly as the article's lines are written, so no omission receipt of any kind accompanies this package. Citations: the retained excerpts carry no citation command at all, so no bibliography entry of the article is transcribed and no reference is redistributed. Reference adaptation: none was needed and none was made; every reference inside the delivered unit resolves inside it, so no unresolved reference or citation is produced. Printed slips kept literal: no printed word, symbol, hypothesis or number of the counted statement or of its proof was altered. Layout and class: the article's own class is \texttt{article} with packages including t1enc, inputenc, babel, amsmath, amsthm, mathrsfs, amsfonts, latexsym, graphicx, float, xcolor, algorithm, algorithmic, enumerate; the wrapper is the lane's pinned \texttt{amsart} editorial wrapper at 10pt on A4, which restates the theorem-like environments the retained text needs and adds the article's own macro definitions that the retained text needs (none), with extra wrapper packages (bm, bbm, dsfont, xspace, cancel, mathtools). The delivered numbering is the unit's own. Assets: no retained span contains a figure environment, a graphics-inclusion command, a PDF-page inclusion, a table or a TikZ picture; no image file is copied into this package, the package declares no asset dependency and no image file is redistributed with it. Licence: version arXiv:2507.04257v1 of this article is distributed under the Creative Commons Attribution 4.0 International licence, as recorded in the arXiv licence metadata for this exact version and shown on the abstract page https://arxiv.org/abs/2507.04257v1. A full rights notice is delivered at licenses/arxiv-2507.04257-CC-BY-4.0.txt. Characters: every retained excerpt is pure ASCII, so the delivered engine, XeLaTeX with its default Latin Modern text font, reports no missing character.
\begin{lemma}\label{lem3.1}
      For every $v\in S'$, % and each $H\in\mathbb{H}$, 
    we have $d_{S''}(v)\leq \max_{H\in \mathbb{H}}\{|H|\}-2$; for every $v\in S''$, we have $d_{S''}(v)\leq\alpha_{\mathbb{H}}$.
\end{lemma}

\begin{lemma}\label{lem3.2}
    For every vertex  $v\in S$ we have $x_v\leq\frac{x_{u^*}}{20C_\mathbb{H}}$.
\end{lemma}

\begin{claim}\label{lem3.4}
If $e(L,S')\neq |L||S'|-1$, then $\rho(G_1)>\rho^*$.
\end{claim}

\begin{claim}\label{lem3.5}
 $\rho(G_2)>\rho^*$. 
\end{claim}

\begin{proof}[\bf Proof of Claim \ref{lem3.5}]
    Let ${\bf y}=(y_1,y_2,\ldots,y_n)^{T}$ be a non-negative unit eigenvector of 
$A(G_1)$ with respect to $\rho(G_1)$, and let ${\bf z}=(z_1,z_2,\ldots, z_n)^T$ be a non-negative unit eigenvector of $A(G_2)$ with respect to $\rho(G_2)$. Let %$y_{u^*}:=\max_{w\in V}y_w$, 
    $y_{S''}^*:=\max_{w\in S''}y_w$,  $\sigma_L:=\sum_{u\in L}y_u$, and let $z_{S''}^*=\max_{u\in S''}z_u$. 
We proceed by  considering the following two cases.

{\bf Case 1.} $e_{G^*}(L,S')\neq |L||S'|-1$.
    
    In view of Claim~\ref{lem3.4}, we have $\rho(G_1)>\rho^*$, so it suffices to prove $\rho(G_2)>\rho(G_1)$. Notice that $\rho(G_1)y_{S''}^*\leq \alpha_{\mathbb{H}}y_{S''}^*+\sigma_L$, this implies 
    \begin{align}\label{ys}
        y_{S''}^*\leq \frac{\sigma_L}{\rho(G_1)-\alpha_{\mathbb{H}}}.
    \end{align}
        %\rho(G_1)y_w=\underset{u\in N_{G_1}(w)}{\sum}y_u=\underset{u\in N_{L\cup S''}(w)}{\sum}y_u+\underset{u\in N_{S'}(w)}{\sum}y_u\leq C_Hy_{u^*}+\frac{\gamma_Hn}{100C_H\rho(G_1)}y_{u^*}%\leq C_Hy_{u^*}+\frac{3\gamma_Hn}{100C_H\rho(G_1)}y_{u^*}.
  

By the construction of
$G_1$, we have $|E(G_1)|\leq |E(G^*)|+|L||S'|\leq C_\mathbb{H} n+\frac{\gamma_\mathbb{H}^2n}{(10C_\mathbb{H})^2}$. Together with the well-known fact that 
$\rho(G_1)\leq \sqrt{2|E(G_1)|}$, we have 
$$
\sqrt{2\left(C_\mathbb{H} n+\frac{\gamma_\mathbb{H}^2n}{(10C_\mathbb{H})^2}\right)}\rho(G_1)y_{v_{j_k}}\geq \rho(G_1)^2y_{v_{j_k}}\geq \sum_{u\in S\setminus \{u_k\}}\rho(G_1)y_u\geq (|S|-1)\sigma_L=(n-\gamma_\mathbb{H}-1)\sigma_L.
$$
It follows that 
\begin{align}\notag
    y_{v_{j_k}}&\geq \frac{(n-\gamma_\mathbb{H}-1)\sigma_L}{\rho(G_1)\sqrt{2\left(C_\mathbb{H} n+\frac{\gamma_\mathbb{H}^2n}{(10C_\mathbb{H})^2}\right)}}\\\notag
    &\geq \frac{\sigma_L}{\rho(G_1)-\alpha_{\mathbb{H}}}\cdot\frac{n-\gamma_\mathbb{H}-1}{2\sqrt{C_\mathbb{H} n+\frac{\gamma_\mathbb{H}^2n}{(10C_\mathbb{H})^2}}}\\\label{eq:4.1}
    &>\frac{\sigma_L\cdot 4\zeta\alpha_{\mathbb{H}}}{\rho(G_1)-\alpha_{\mathbb{H}}}\geq 4\zeta\alpha_{\mathbb{H}}y_{S''}^*.
\end{align}

    We first assume that $\gamma_\mathbb{H}\geq2$. Clearly, 
    \begin{align}\notag\label{eq:4}
        \rho(G_2)-\rho(G_1)&\geq{\bf y}^T\left(A(G_2)-A(G_1)\right){\bf y}\\&\geq 2y_{v_{j_k}}y_{u_k}-2\sum_{i=1}^{\zeta}\sum_{w\in N_{S''}(w_i)}y_{w_i}y_w.
    \end{align}
We claim that $\sigma_L\geq 2y_{v_{j_k}}$. In fact, if $y_{v_{j_k}}\leq y_{v_{\ell}}$ for some $\ell\neq j_k$, then $\sigma_L\geq 2y_{v_{j_k}}$. If $y_{v_{j_k}}> y_{v_{\ell}}$ for every  $\ell\neq j_k$, then $N_L(v_{j_k})\neq \emptyset$ and  
$$
\rho(G_1)\left(\sum_{v_\ell\in N_L(v_{j_k})}y_{v_\ell}-y_{v_{j_k}}\right)\geq |N_L(v_{j_k})|y_{v_{j_k}}-\sum_{v_\ell\in N_L(v_{j_k})}y_{v_\ell}>0.
$$
Therefore, $\sigma_L\geq 2y_{v_{j_k}}$, as desired. 
Notice that $\rho(G_1)y_{u_k}=\sigma_L-y_{v_{j_k}}$, % and  $\rho(G_1)y_{S''}^*\leq\sigma_L+\alpha_Hy_{S''}^*$ for each $H\in\mathbb{H}$, 
which combined with \eqref{ys}, implies that  
\begin{align*}
y_{S''}^*\leq\frac{\sigma_L}{\rho(G_1)-\alpha_\mathbb{H}}\leq\frac{2(\sigma_L-y_{v_{j_k}})}{0.5\rho(G_1)}=4y_{u_k}.
\end{align*}
   Therefore, %for each $H\in\mathbb{H}$,
    \begin{align*}
        \sum_{i=1}^{\zeta}\sum_{w\in N_{S''}(w_i)}y_{w_i}y_w\leq\zeta\alpha_\mathbb{H}\cdot y_{S''}^*\cdot4y_{u_k}=4\zeta\alpha_\mathbb{H}\cdot y_{S''}^*y_{u_k}.
    \end{align*}
Combining this with \eqref{eq:4.1} and \eqref{eq:4}, we get that $\rho(G_2)-\rho(G_1)>0$, as desired. 
    
    Next, we assume that $\gamma_\mathbb{H}=1$ and $v_1 (=v_{j_k})$ is the unique vertex in $L$.
%Recall that ${\bf z}=(z_1,z_2,\ldots,  z_n)^T$ is a non-negative unit eigenvector of $A(G_2)$ with respect to $\rho(G_2)$.
    Notice that $\rho(G_2)z_{S''}^*\leq z_{v_1}+\alpha_\mathbb{H}z_{S''}^*$ %for each $H\in\mathbb{H}$
    and $\rho(G_2)z_{u_k}=z_{v_1}$. 
    This yields $(\rho(G_2)-\alpha_\mathbb{H})z_{S''}^*\leq \rho(G_2)z_{u_k}$, % for each $H\in\mathbb{H}$ , 
    which implies $z_{S''}^*\leq 2z_{u_k}$. Combining this with \eqref{eq:4.1}, one has 
    %Using the same method in~\eqref{eq2}, we have $z_{S''}^*<\frac{z_{u^*}}{20C_H}$.
    %Thus,
    \begin{align*}
        {\bf y}^T{\bf z}(\rho(G_2)-\rho(G_1))&={\bf y}^T\left(A(G_2)-A(G_1)\right){\bf z}\\
        &=y_{u_k}z_{v_1}+y_{v_1}z_{u_k}-\sum_{i=1}^{\zeta}\sum_{w\in N_{S''}(w_i)}(y_{w_i}z_w+y_wz_{w_i})\\&\geq y_{v_1}z_{u_k}-2\zeta\cdot\alpha_\mathbb{H}\cdot y_{S''}^*z_{S''}^*\\&>0.%\geq y_{v_1}z_{u_k}\left(1-\frac{1}{20C_H}\cdot8|H|^2\alpha_H\right)\\&>0
    \end{align*}
That is, $\rho(G_2)>\rho(G_1)$, as desired. 

{\bf Case 2.} $e_{G^*}(L,S')=|L||S'|-1$.

In this case, $S'=\{u_1\}$ and $d_L(u_1)=\gamma_\mathbb{H}-1$. 
    %Recall that ${\bf x}=(x_1,x_2,\ldots, x_n)^T$ is a non-negative unit eigenvector of $G^*$ corresponding to $\rho^*$ and ${\bf z}=(z_1,z_2\ldots z_n)^T$ is a non-negative unit eigenvector of $A(G_2)$ with respect to $\rho(G_2)$. Define $z_{S''}^*:=\max_{w\in S''}z_w$. 
    Notice that $\rho(G_2)z_{u_1}=\sum_{w\in L}z_w$ and $\rho(G_2)z_{S''}^*\leq\sum_{w\in L}z_w+\alpha_\mathbb{H}z_{S''}^*$, % for each $H\in\mathbb{H}$, 
    which implies $z_{S''}^*\leq 2z_{u_1}$. Recall that $x_{v_{j_1}} \geq \left(1-\frac{1}{2(10C_\mathbb{H})^{2}}\right)x_{u^{*}}$.  Together with Lemmas~\ref{lem3.1}-\ref{lem3.2} and the choice of $C_\mathbb{H}$, one has
    \begin{align*}
        {\bf x}^T{\bf z}(\rho(G_2)-\rho^*)&={\bf x}^T(A(G_2)-A(G^*)){\bf z} \\&\geq x_{u_1}z_{v_{j_1}}+x_{v_{j_1}}z_{u_1}-\sum_{i=1}^{\zeta}\sum_{w\in N_{S''}(w_i)}(x_{w_i}z_w+x_wz_{w_i})-\sum_{w\in N_{S''}(u_1)}(x_{u_1}z_w+x_wz_{u_1})\\ &\geq x_{v_{j_1}}z_{u_1}-2\zeta\cdot \alpha_\mathbb{H}\cdot x_{S''}^*z_{S''}^*-\zeta \cdot x_{S''}^*z_{S''}^*\\ &\geq \frac{1}{2}x_u^*z_{u_1}-3\zeta\cdot  \alpha_\mathbb{H}\cdot\frac{x_u^*}{20C_\mathbb{H}}\cdot2z_{u_1}\\ &>0.
    \end{align*}
That is, $\rho(G_2)>\rho^*$, as desired. 
\end{proof}

\end{document}
